\chapter{La mort, le repos, l’intercession et l’espérance}\label{ch:mort_esperance}

\section{Le repos et le sommeil}

Le vocabulaire de la mort est dominé par le repos. La formule \emph{in pace} est attestée dans \st{ip_tot} textes funéraires, dont \st{ip_sansrest} sans restitution (chapitre~\ref{ch:formulaire}). Les verbes \emph{quiescere} et \emph{requiescere} sont lus dans 29 textes et restitués dans 8 autres ; ils se rencontrent au présent (\emph{hic quiescit}, \emph{requiescit in pace}), sous la forme \emph{quiescet} (\nnums{42}, 59 et 102), qui peut noter un futur ou, plus probablement, le présent avec la confusion de \emph{e} et de \emph{i} (chapitre~\ref{ch:langue}), et à la deuxième personne, dans l’adresse au défunt (\emph{quiescis}, \nnums{73} et 231 ; peut-être \emph{q(uiescas)}, simple initiale développée par l’éditeur, au \nnum{446}). \emph{Pausare} n’est lu qu’une fois (\emph{pausavit}, \nnum{458}).

L’image du sommeil est plus rare : \emph{dormitio} pour désigner la mort (\nnum{134}, où le mot ouvre deux des commémorations ; restitué au \nnum{14}), \emph{dormit} (\nnums{73} et 400 ; restitué au \nnum{464}), \emph{in somn\{i\}o pacis} (\nnum{400}, lu ; restitué aux \nnums{257} et 418). Le sommeil peut s’accorder à une espérance de réveil ; le mot seul ne la formule pas.

D’autres expressions désignent la mort comme un départ : \emph{recessit}, \emph{decessit}, \emph{[d]isces[sit de ha]c luce} (\nnum{309}, en partie restitué), \emph{absolutus de corpore} (\nnum{211}, en partie restitué), \emph{praecessit ad pacem} (\nnum{275}), \emph{exivit} (\nnum{271}), \emph{fatum fecit} (\nnum{295}). La tombe est une \emph{domus aeterna} (\nnum{126} ; en partie restitué aux \nnums{73} et 309) : l’expression, héritée de l’épigraphie funéraire païenne, est employée sans difficulté dans des textes qui portent \emph{in pace}.

\section{Persistances du langage traditionnel}

Le vocabulaire chrétien du repos coexiste avec des formules traditionnelles. Le \nnum{72} (362) s’ouvre sur \emph{D(is) M(anibus) s(acrum)} et se poursuit par \emph{requievit in pace} : c’est le seul texte de l’ensemble qui porte la dédicace aux Mânes. Les poèmes empruntent à la tradition classique : \emph{nescit fama mori} pour le juriste Caianus (\nnum{22}, 348) ; \emph{nec curant carmina manes} dans l’épitaphe de Felicitas (\nnum{470}, 403), qui se termine pourtant par \emph{dep(osita) in p(ace)} ; \emph{si fas} au \nnum{440} (400). Ces emprunts relèvent d’un langage poétique partagé ; ils ne signalent pas une hésitation religieuse des commanditaires.

\section{L’au-delà : ciel, lumière, saints}

Une dizaine de textes, presque tous poétiques, évoquent le sort du défunt au-delà de la mort. Les formulations sont diverses, et plusieurs reposent en partie sur des restitutions.

\begin{itemize}
\item \emph{La lumière.} L’esprit de Severianus « a été accueilli dans la lumière du Seigneur » (\nnum{400}, lu) ; Eutychius, « croyant au Christ, possède la récompense de la lumière » (\emph{praemia lucis habet}, \nnum{352}) ; Proiecta, « désirant monter vers la lumière éthérée du ciel » (\emph{aetheriam cupiens caeli conscendere lucem}, \nnum{244}, épitaphe due à l’évêque Damase, lue).
\item \emph{Le ciel.} La défunte du \nnum{382} « possède les royaumes célestes » (\emph{celestia regna t[enet]}, verbe en partie restitué) ; au \nnum{433}, l’enfant « se réjouit parmi les astres » et « est allé aux astres », mais les mots \emph{astris} et \emph{astra} appartiennent entièrement aux fins de vers restituées, et l’EDB souligne de nombreuses lettres du reste du poème, lues par les éditeurs anciens et aujourd’hui perdues.
\item \emph{La vie éternelle.} L’épitaphe de Maximilla (\nnum{306}) développe l’argument le plus construit du corpus : la mort a emporté tous les biens de la vie, « mais il vit mieux, celui qui espère la vie éternelle du Christ » ; « alors la mort elle-même pourra mourir, et la vie sera éternelle pour tous les temps » (\emph{tunc poterit mors ipsa mori cum tempore toto vita perennis erit}). Le \nnum{231} parle de « gloire éternelle » (\emph{gloria perennis}) et d’« espérance future » (\emph{spes futura}) pour un enfant mort à seize mois.
\item \emph{Les saints.} Quatre textes associent le défunt aux saints : au \nnum{73}, l’esprit de la défunte « vit avec » ou « est associé aux » saints, selon la restitution du verbe (\emph{[vivit cum]} dans l’EDCS, \emph{[est sociatu]s} dans l’EDB), \emph{pro meritis et opera tanta} ; au \nnum{232}, la défunte a obtenu une sépulture « au seuil des saints », \emph{[intra l]imina sanctorum}, « ce que beaucoup désirent et que peu obtiennent » (formule en partie restituée) ; au \nnum{330}, le sigle \emph{p. t. c. s.} est développé \emph{p(ax) t(ibi) c(um) s(anctis)} ; au \nnum{348}, Alexandros est invité à « vivre avec les saints » (\emph{ζέσας μετὰ τῶν ἁγίων}, si \emph{ζέσας} est bien une graphie de l’optatif \emph{ζήσαις}).
\end{itemize}

Ces expressions relèvent pour la plupart d’un langage commun à la poésie funéraire de l’époque. Elles montrent que l’au-delà est pensé comme lumière, ciel, repos et société des saints ; elles ne permettent pas, à elles seules, de reconstituer une doctrine des fins dernières.

\section{La résurrection : aucune mention assurée}

Aucune épitaphe latine du corpus ne nomme la résurrection. Le seul texte qui pourrait la mentionner est l’imprécation grecque du \nnum{447}, où l’EDCS restitue \emph{ἡμαίραν τῆς [ἀναστάσεως(?)]}, avec un point d’interrogation, et où l’EDB restitue \emph{τὴν φοβερὰν ἡμαίραν τῆς [κρίσεως]}, « le jour redoutable du jugement » (chapitre~\ref{ch:grec}). La version précédente de ce livre comptait une mention de la résurrection ; elle reposait sur une seule des deux restitutions publiées. Le texte transmis atteste seulement la menace d’un « jour » eschatologique.

L’absence du mot n’est pas l’absence de la croyance. L’image du sommeil, l’accueil de l’esprit dans la lumière, le corps confié à la terre « comme un gage » (\emph{mandasti corpus terrae pignusque}, \nnum{73}, en partie restitué) peuvent s’accorder à l’espérance de la résurrection. Mais aucun texte du corpus ne l’énonce.

\section{L’intercession}

Deux textes adressent au défunt une demande de prière. Au \nnum{214} (380), la défunte, \emph{ancilla dei}, est priée d’intercéder pour l’enfant qu’elle laisse : \emph{pro hunc unum ora subolem quem superstitem reliquisti}, « prie pour ce seul enfant que tu as laissé survivant ». La phrase est lue dans les deux bases. Au \nnum{457}, un fragment de date incertaine (l’année 401 est celle de la base), les mots \emph{[ora]tionibus t[uis]} et \emph{[pro fil]io et fr[atribus]} sont en grande partie restitués. Le \nnum{214} est donc le seul témoignage assuré du corpus. Il suppose la conviction que le défunt, auprès de Dieu, peut prier pour les vivants.
